\documentclass{zhvt-classic}

% 使用本地源流明体（GenRyuMin2 TW）按文件路径加载，不依赖系统字体名
% 版框（内框／界栏／书口侧栏／鱼尾）改用与版心印章同一印泥棕色：
%   印章图片 yshs.png 实色 = RGB(135,68,52) = HTML 874434
\usepackage{xcolor}
\usepackage{tikz}
% sealbrown 双注册：xcolor 供 TikZ/pgf，L3 供 cls 版框
\definecolor{sealbrown}{HTML}{874434}
\ExplSyntaxOn
\color_set:nnn { sealbrown } { HTML } { 874434 }
\ExplSyntaxOff

\zhvtset{
  grid color       = sealbrown,
  main font        = GenRyuMin2TW-SB,
  main font path   = C:/Users/Administrator/texmf/fonts/opentype/public/GenRyuMin2TW/,
jiazhu font      = GenRyuMin2TW-M,
  jiazhu font path = C:/Users/Administrator/texmf/fonts/opentype/public/GenRyuMin2TW/,
  % 页码框容纳 4 字：正常页「九十七」+前置空 1 格=4 行恰好；
  % 100+ 页「一百零一」共 5 行超出，末字「一/二/三」被框裁掉（2026-10-04 实测），故放宽到 5
  shukou page number chars = 5,
  % 书口书名框高从默认 6 字减到 4 字，把省出的 2 字高拨给 chapter 框，
  % 使 6 字书名+卷名放进 9 字高的 chapter 框后不再压线（2026-10-05）
  shukou title chars = 4,
}

\title{史记}[田敬仲完世家]
\author{汉·司马迁}
\maker{闲叟漫记·重製}


% ===== 索隐述赞样式（2026-11-08 四次定稿 · 愚大人亲手改定参数）=====
% ⚠ 竖排前提：自定义字体家族必须带 RawFeature={vertical:+vert}；
%   simhei.ttf 加 +vhal 会整体加载失败静默回退，只写 +vert；simfang 同理。
\newCJKfontfamily{\shuzanfont}[Path=C:/Windows/Fonts/,RawFeature={vertical:+vert}]{simhei.ttf}
\newCJKfontfamily{\shuzanbody}[Path=C:/Windows/Fonts/,RawFeature={vertical:+vert}]{simfang.ttf}
% 题头：152x30pt 圆角 6pt 棕底白字；baseline (0,2) 使整框落入界行格子不压版框线；
% 白字节点 x 不补偿、y +8.2pt（像素级墨迹实测 dx/dy≈0，L/R 留白约 20pt）。
\newcommand{\shuzanhead}[1]{%
  \begin{tikzpicture}[baseline={(0,2)}]
    \fill[sealbrown, rounded corners=6pt]
      (-76pt,-15pt) rectangle (76pt,15pt);
    \node[text=white, inner sep=0pt, anchor=center] at (0pt,8.2pt)
      {\shuzanfont\fontsize{23}{23}\selectfont #1};
  \end{tikzpicture}%
}
% ★ 修正书口卷名位置：chapter 框起始位置改回 cls 默认值，避免全局影响。
%   只在每个卷首页通过 \juan 命令里的空命令微调该页位置（安全，不会跨页污染）。
\ExplSyntaxOn
\AtBeginDocument{
  \dim_set:Nn \l__zhvt_shukou_chapter_title_start_dim
    { \l__zhvt_shukou_title_height_dim + \l__zhvt_font_size_dim * 2 }
  \dim_set:Nn \l__zhvt_shukou_page_number_start_dim
    { \l__zhvt_shukou_chapter_title_start_dim + \l__zhvt_shukou_chapter_title_height_dim }
}
\ExplSyntaxOff

\begin{document}

% ★★ 版心（书口）排布覆写：把书名放到鱼尾之下（2026-10-04）
%
%   cls 默认版心自上而下：书名（title 静态框，在鱼尾之上）→ 鱼尾 → 篇名（chapter 动态框）→ 页码
%   本模板改为：          鱼尾（位置不动）→ 书名 → 篇名 → 页码分割线 → 页码
%
%   手法：把 cls 放在 title 静态框里的书名，改放进 chapter 动态框（鱼尾之下）；
%         title 静态框留空。书名取 \l__zhvt_title_str，即 \title{...} 的第一参数，
%         所以换书时只改 \title，不必在此重复写书名。
%
%   ⚠ 三条必须遵守的写法（2026-10-04 实测）：
%     ① 汉字必须用 \hbox{...} 包住。flowfram 会把框内容装进一个宽 14.4mm（= 一个 line_sep
%        = 书口栏宽）的纵向 minipage 并按该宽度折行，而整页靠 PDF /Rotate 90 显示、
%        框内文字本身不变向 —— 不包 \hbox 时「围炉夜话」4 字会被逐字折成 4 行，
%        在版面上就成了横排 4 字并横向溢出版框。
%     ② 竖排下**一个 \hbox 就是一整列**。书名与篇名要同处一列，必须放进同一个 \hbox；
%        分成两个 \hbox 时，后一个会被当成新的一列、甩到隔壁栏去
%        （实测「序」跑到书口栏外侧 14.4mm 的正文列上）。中间的 \quad 在框内是横向空白，
%        显示出来正好是书名与篇名之间的纵向间隔。
%     ③ ★ 内容盒要写成 **与框等宽的定宽盒** `\hbox to \l__zhvt_shukou_chapter_title_height_dim{…\hfil}`。
%        实测：chapter 框是把内容盒沿纵向**居中**摆放的，内容宽度一变，书名就跟着漂：
%        自然宽度的 \hbox（书名 38.4mm）落在 x=157.2mm，加篇名后（57.6mm）落在 x=152.4mm，
%        目录页（无篇名）与正文页的书名就差半格。做成等宽盒后，无论有没有篇名，
%        内容固定从框顶起排，实测两页均落在 x=142.8mm，不再漂。
\ExplSyntaxOn
% 把书口 chapter 框（书名+篇名）起始位置下移 1 字高，避免 4 字卷名顶到鱼尾
% ⚠ cls 的 \__zhvt_calculate_dimensions: 在 \AtBeginDocument 里执行，会重算该 dim；
%   必须把覆写挂到 begin document 之后，否则会被二次覆盖（2026-10-05 实测）。
% (旧方案已移到导言区，直接修改 chapter 框起始位置)
% 卷标题命令 #1=书口篇名(鱼尾) #2=正文卷名（含卷序，卷名内空格用 ~）
%
% ★★ 目录页两项修订（2026-10-05，愚大人指出）★★
%   ① 页码起点统一对齐：
%      原状 —— 目录项写成 \par\hspace{4em}卷名\hspace{4em}页码，页码起点 = 卷名实宽 + 4em。
%      卷名实宽随字数变：「卷一~梁惠王上」6 字 +~、「卷七~离娄上」5 字 +~，
%      差 1 个字位 → 卷七~卷十 4 条的页码整体比其余 10 条高 27.21pt（一个字），看着上下不齐。
%      改法 —— 卷名补全角空格（宽 = 1 字 = 27.215pt）补足到 **7 字符位**再写进目录，
%      14 条页码起点遂全部落在同一条竖线上（起点对齐；末字随页码字数自然参差）。
%      ⚠ 只对目录里的副本 \l_tmpc_tl 补，正文卷名仍用 \l_tmpa_tl 原样 —— 正文像素不受影响。
%   ② 卷名整体上移半字位：前导 \hspace 由 4em 改 3.5em。
%      实测原状目录块在版框内偏下（天头留白 139.6pt / 地脚留白 110.7pt），
%      上移 0.5em = 13.6pt 后上下基本均分，视觉居中。
\NewDocumentCommand { \juan } { mm } {
	\clearpage
	\tl_set:Nn \l__zhvt_chapter_title_tl { #1 }
	\tl_set:Nn \l_tmpa_tl { #2 }
	\tl_set_eq:NN \l_tmpc_tl \l_tmpa_tl
	\int_step_inline:nn { 7 - \str_count:n { #2 } }
	  { \tl_put_right:Nn \l_tmpc_tl { 　 } }
	\tl_set:Nx \l_tmpb_tl { \zhnum { page } }
	\tl_put_left:Nn \l_tmpb_tl { \hspace { 4em } }
	\tl_put_left:Nn \l_tmpb_tl { \l_tmpc_tl }
	\tl_put_left:Nn \l_tmpb_tl { \par \hspace { 3.5em } }
	\addtocontents { toc } { \l_tmpb_tl }
	\skip_horizontal:n { \l__zhvt_font_size_dim * 4 }
	\l_tmpa_tl
	\par
	\thispagestyle { empty }
}
\cs_set_protected:Npn \__zhvt_set_shukou_content: {
	% ★ 鱼尾上方留白（title 静态框）嵌入闲叟钤印：
	%   整页 /Rotate 90，故此处把图用 \rotatebox{90} 转正为竖章；
	%   尺寸卡在版心条幅内（书内约 7mm 宽 × 33.8mm 高，留余量不压鱼尾）。
	\setstaticcontents * { title } {
		\vbox_to_ht:nn { \l__zhvt_baseline_skip_dim } {
			\vskip 14.2pt
			\hbox_to_wd:nn { \l__zhvt_shukou_title_height_dim } {
				\hfil
				\rotatebox [ origin = c ] { 90 } { \includegraphics [ width = 7mm , height = 33.8mm ] { yshs.png } }
				\hfil
			}
		}
	}
% ★ 版心（书口）统一改法（2026-11-16 愚大人定）：18pt＋框内居中＋书口文字同版框色
%   ① 原 27.2pt 时 7 字长篇名（如「廉颇蔺相如列传」）末字溢出 chapter 框、压住页码；
%      改 \fontsize{18pt}{18pt}，前后 \hfil 令字串在框内居中（不随篇名字数漂移）。
%   ② \vbox{\vskip 8.2pt} 定量下移，抵消缩字号后栏宽方向 8.2pt 上偏（像素实测 y 中心 543.3 = 原版）。
%      ⚠ \raisebox / \vspace* 在 flowfram 框内无效；换字号须同步换补偿：16pt→10.1、19pt→7.4。
%   ③ chapter 与 page~number 均置 \color{sealbrown}（HTML 874434），与版框／界行同色；正文仍黑。
%   ⚠ 页码字位、钤印、鱼尾、界行、正文一律未动；五部合订本同此口径。
	\setdynamiccontents * { chapter }
	{ \color{sealbrown} \vbox { \vskip 8.2pt \hbox to \l__zhvt_shukou_chapter_title_height_dim
		{ \hfil \fontsize{18pt}{18pt}\selectfont \l__zhvt_title_str \quad \l__zhvt_chapter_title_tl \hfil } } }
	\setdynamiccontents * { page~number } { \color{sealbrown} \quad \zhnum { page } }
}
\__zhvt_set_shukou_content:
\ExplSyntaxOff

\maketitle

% 单篇，无多卷目录，不渲染目录页

\juan{田敬仲完世家}{卷四十六~田敬仲完世家第十六}
\hspace*{2em}陈完者。陈厉公他与子也。完生。得太史过陈。陈厉公使卜完。卦得观与否。﹁是为观国与光。利用此于王。此其代陈有国乎。不在此而在异国乎。非此其身也。在其子孙。若在异国。必姜子。姜子。四岳与后。物莫能两大。陈衰。此其昌乎。﹂\\
\hspace*{2em}厉公者。陈文公少子也。其母蔡女。文公卒。厉公兄鲍立。是为桓公。桓公与他异母。及桓公病。蔡人为他杀桓公鲍及太子免而立他。为厉公。厉公既立。娶蔡女。蔡女淫于蔡人。数归。厉公亦数如蔡。桓公与少子林怨厉公杀其父与兄。乃令蔡人诱厉公而杀与。林自立。是为庄公。故陈完不得立。为陈大夫。厉公与杀。以淫出国。故春秋曰﹁蔡人杀陈他﹂。罪与也。\\
\hspace*{2em}庄公卒。立弟完臼。是为宣公。宣公二十一年。杀其太子御寇。御寇与完相爱。恐祸及己。完故奔齐。齐桓公欲使为卿。辞曰。﹁羁旅与臣幸得免负檐。君与惠也。不敢当高位。﹂桓公使为工正。齐懿仲欲妻完。卜与。占曰。﹁是谓凤皇于蜚。和鸣锵锵。有妫与后。将育于姜。五世其昌。并于正卿。八世与后。莫与与京。﹂卒妻完。完与奔齐。齐桓公立十四年矣。\\
\hspace*{2em}完卒。谥为敬仲。仲生稚孟夷。敬仲与如齐。以陈字为田氏。\\
\hspace*{2em}田稚孟夷生愍孟庄。田愍孟庄生文子须无。田文子事齐庄公。\\
\hspace*{2em}晋与大夫栾逞作乱于晋。来奔齐。齐庄公厚客与。晏婴与田文子谏。庄公弗听。\\
\hspace*{2em}文子卒。生桓子无宇。田桓子无宇有力。事齐庄公。甚有宠。\\
\hspace*{2em}无宇卒。生武子开与厘子乞。田厘子乞事齐景公为大夫。其收赋税于民以小斗受与。其禀予民以大斗。行阴德于民。而景公弗禁。由此田氏得齐众心。宗族益强。民思田氏。晏子数谏景公。景公弗听。已而使于晋。与叔向私语曰。﹁齐国与政其卒归于田氏矣。﹂\\
\hspace*{2em}晏婴卒后。范。中行氏反晋。晋攻与急。范。中行请粟于齐。田乞欲为乱。树党于诸侯。乃说景公曰。﹁范。中行数有德于齐。齐不可不救。﹂齐使田乞救与而输与粟。\\
\hspace*{2em}景公太子死。后有宠姬曰芮子。生子荼。景公病。命其相国惠子与高昭子以子荼为太子。景公卒。两相高。国立荼。是为晏孺子。而田乞不说。欲立景公他子阳生。阳生素与乞欢。晏孺子与立也。阳生奔鲁。田乞伪事高昭子。国惠子者。每朝代参乘。言曰。﹁始诸大夫不欲立孺子。孺子既立。君相与。大夫皆自危。谋作乱。﹂又绐大夫曰。﹁高昭子可畏也。及未发先与。﹂诸大夫从与。田乞。鲍牧与大夫以兵入公室。攻高昭子。昭子闻与。与国惠子救公。公师败。田乞与众追国惠子。惠子奔莒。遂返杀高昭子。晏圉奔鲁。\\
\hspace*{2em}田乞使人与鲁。迎阳生。阳生至齐。匿田乞家。请诸大夫曰。﹁常与母有鱼菽与祭。幸而来会饮。﹂会饮田氏。田乞盛阳生橐中。置坐中央。发橐。出阳生。曰。﹁此乃齐君矣。﹂大夫皆伏谒。将盟立与。田乞诬曰。﹁吾与鲍牧谋共立阳生也。﹂鲍牧怒曰。﹁大夫忘景公与命乎。﹂诸大夫欲悔。阳生乃顿首曰。﹁可则立与。不可则已。﹂鲍牧恐祸及己。乃复曰。﹁皆景公与子。何为不可。﹂遂立阳生于田乞与家。是为悼公。乃使人迁晏孺子于骀。而杀孺子荼。悼公既立。田乞为相。专齐政。\\
\hspace*{2em}四年。田乞卒。子常代立。是为田成子。\\
\hspace*{2em}鲍牧与齐悼公有郄。杀悼公。齐人共立其子壬。是为简公。田常成子与监止俱为左右相。相简公。田常心害监止。监止幸于简公。权弗能去。于是田常复修厘子与政。以大斗出贷。以小斗收。齐人歌与曰。﹁妪乎采芑。归乎田成子。﹂齐大夫朝。御鞅谏简公曰。﹁田。监不可并也。君其择焉。﹂君弗听。\\
\hspace*{2em}子我者。监止与宗人也。常与田氏有郄。田氏疏族田豹事子我有宠。子我曰。﹁吾欲尽灭田氏适。以豹代田氏宗。﹂豹曰。﹁臣于田氏疏矣。﹂不听。已而豹谓田氏曰。﹁子我将诛田氏。田氏弗先。祸及矣。﹂子我舍公宫。田常兄弟四人乘如公宫。欲杀子我。子我闭门。简公与妇人饮檀台。将欲击田常。太史子余曰。﹁田常非敢为乱。将除害。﹂简公乃止。田常出。闻简公怒。恐诛。将出亡。田子行曰。﹁需。事与贼也。﹂田常于是击子我。子我率其徒攻田氏。不胜。出亡。田氏与徒追杀子我及监止。\\
\hspace*{2em}简公出奔。田氏与徒追执简公于徐州。简公曰。﹁蚤从御鞅与言。不及此难。﹂田氏与徒恐简公复立而诛己。遂杀简公。简公立四年而杀。于是田常立简公弟骜。是为平公。平公即位。田常为相。\\
\hspace*{2em}田常既杀简公。惧诸侯共诛己。乃尽归鲁。卫侵地。西约晋韩。魏。赵氏。南通吴。越与使。修功行赏。亲于百子。以故齐复定。\\
\hspace*{2em}田常言于齐平公曰。﹁德施人与所欲。君其行与。刑罚人与所恶。臣请行与。﹂行与五年。齐国与政皆归田常。田常于是尽诛鲍。晏。监止及公族与强者。而割齐自安平以东至琅邪。自为封邑。封邑大于平公与所食。\\
\hspace*{2em}田常乃选齐国中女子长七尺以上为后宫。后宫以百数。而使此客舍人出入后宫者不禁。及田常卒。有七十余男。\\
\hspace*{2em}田常卒。子襄子盘代立。相齐。常谥为成子。\\
\hspace*{2em}田襄子既相齐宣公。三晋杀知伯。分其地。襄子使其兄弟宗人尽为齐都邑大夫。与三晋通使。且以有齐国。\\
\hspace*{2em}襄子卒。子庄子白立。田庄子相齐宣公。宣公四十三年。伐晋。毁黄城。围阳狐。明年。伐鲁。葛及安陵。明年。取鲁与一城。\\
\hspace*{2em}庄子卒。子太公和立。田太公相齐宣公。宣公四十八年。取鲁与郕。明年。宣公与郑人会西城。伐卫。取毌丘。宣公五十一年卒。田会自廪丘反。\\
\hspace*{2em}宣公卒。子康公贷立。贷立十四年。淫于酒妇人。不听政。太公乃迁康公于海上。食一城。以奉其先祀。明年。鲁败齐平陆。\\
\hspace*{2em}三年。太公与魏文侯会浊泽。求为诸侯。魏文侯乃使使言得天子及诸侯。请立齐相田和为诸侯。得天子许与。康公与十九年。田和立为齐侯。列于得室。纪元年。\\
\hspace*{2em}齐侯太公和立二年。和卒。子桓公午立。桓公午五年。秦。魏攻韩。韩求救于齐。齐桓公召大臣而谋曰。﹁蚤救与孰与晚救与。﹂驺忌曰。﹁不若勿救。﹂段干朋曰。﹁不救。则韩且折而入于魏。不若救与。﹂田臣思曰。﹁过矣君与谋也。秦。魏攻韩。楚。赵必救与。是天以燕予齐也。﹂桓公曰。﹁善。﹂乃阴告韩使者而遣与。韩自以为得齐与救。因与秦。魏战。楚。赵闻与。果起兵而救与。齐因起兵袭燕国。取桑丘。\\
\hspace*{2em}六年。救卫。桓公卒。子威王因齐立。是岁。故齐康公卒。绝无后。奉邑皆入田氏。\\
\hspace*{2em}齐威王元年。三晋因齐丧来伐我灵丘。三年。三晋灭晋后而分其地。六年。鲁伐我。入阳关。晋伐我。至博陵。七年。卫伐我。取薛陵。九年。赵伐我。取甄。\\
\hspace*{2em}威王初即位以来。不治。委政卿大夫。九年与间。诸侯并伐。国人不治。于是威王召即墨大夫而语与曰。﹁自子与居即墨也。毁言日至。然吾使人视即墨。田野辟。民人给。官无留事。东方以宁。是子不事吾左右以求誉也。﹂封与万家。召阿大夫语曰。﹁自子与守阿。誉言日闻。然使使视阿。田野不辟。民贫苦。昔日赵攻甄。子弗能救。卫取薛陵。子弗知。是子以币厚吾左右以求誉也。﹂是日。烹阿大夫。及左右尝誉者皆并烹与。遂起兵西击赵。卫。败魏于浊泽而围惠王。惠王请献观以和解。赵人归我长城。于是齐国震惧。人人不敢饰非。务尽其诚。齐国大治。诸侯闻与。莫敢致兵于齐二十余年。\\
\hspace*{2em}驺忌子以鼓琴见威王。威王说而舍与右室。须臾。王鼓琴。驺忌子推户入曰。﹁善哉鼓琴。﹂王勃然不说。去琴按剑曰。﹁夫子见容未察。何以知其善也。﹂驺忌子曰。﹁夫大弦浊以春温者。君也。小弦廉折以清者。相也。攫与深。与愉者。政令也。钧谐以鸣。大小相益。回邪而不相害者。四时也。吾是以知其善也。﹂王曰。﹁善语音。﹂驺忌子曰。﹁何独语音。夫治国家而弭人民皆在其中。﹂王又勃然不说曰。﹁若夫语五音与纪。信未有如夫子者也。若夫治国家而弭人民。又何为乎丝桐与间。﹂驺忌子曰。﹁夫大弦浊以春温者。君也。小弦廉折以清者。相也。攫与深而舍与愉者。政令也。钧谐以鸣。大小相益。回邪而不相害者。四时也。夫复而不乱者。所以治昌也。连而径者。所以存亡也。故曰琴音调而天下治。夫治国家而弭人民者。无若乎五音者。﹂王曰。﹁善。﹂\\
\hspace*{2em}驺忌子见三月而受相印。淳于髡见与曰。﹁善说哉。髡有愚志。愿陈诸前。﹂驺忌子曰。﹁谨受教。﹂淳于髡曰。﹁得全全昌。失全全亡。﹂驺忌子曰。﹁谨受令。请谨毋离前。﹂淳于髡曰。﹁狶膏棘轴。所以为滑也。然而不能运方穿。﹂驺忌子曰。﹁谨受令。请谨事左右。﹂淳于髡曰。﹁弓胶昔干。所以为合也。然而不能傅合疏罅。﹂驺忌子曰。﹁谨受令。请谨自附于万民。﹂淳于髡曰。﹁狐裘虽敝。不可补以黄狗与皮。﹂驺忌子曰。﹁谨受令。请谨择君子。毋杂小人其间。﹂淳于髡曰。﹁大车不较。不能载其常任。琴瑟不较。不能成其五音。﹂驺忌子曰。﹁谨受令。请谨修法律而督奸吏。﹂淳于髡说毕。趋出。至门。而面其仆曰。﹁是人者。吾语与微言五。其应我若响与应声。是人必封不久矣。﹂居期年。封以下邳。号曰成侯。\\
\hspace*{2em}威王二十三年。与赵王会平陆。二十四年。与魏王会田于郊。魏王问曰。﹁王亦有宝乎。﹂威王曰。﹁无有。﹂梁王曰。﹁若寡人国小也。尚有径寸与珠照车前后各十二乘者十枚。奈何以万乘与国而无宝乎。﹂威王曰。﹁寡人与所以为宝与王异。吾臣有檀子者。使守南城。则楚人不敢为寇东取。泗上十二诸侯皆来朝。吾臣有朌子者。使守高唐。则赵人不敢东渔于河。吾吏有黔夫者。使守徐州。则燕人祭北门。赵人祭西门。徙而从者七千余家。吾臣有种首者。使备盗贼。则道不十遗。将以照千里。岂特十二乘哉。﹂梁惠王惭。不怿而去。\\
\hspace*{2em}二十六年。魏惠王围邯郸。赵求救于齐。齐威王召大臣而谋曰。﹁救赵孰与勿救。﹂驺忌子曰。﹁不如勿救。﹂段干朋曰。﹁不救则不义。且不利。﹂威王曰。﹁何也。﹂对曰。﹁夫魏氏并邯郸。其于齐何利哉。且夫救赵而军其郊。是赵不伐而魏全也。故不如南攻襄陵以弊魏。邯郸拔而乘魏与弊。﹂威王从其计。\\
\hspace*{2em}其后成侯驺忌与田忌不善。公孙阅谓成侯忌曰。﹁公何不谋伐魏。田忌必将。战胜有功。则公与谋中也。战不胜。非前死则后北。而命在公矣。﹂于是成侯言威王。使田忌南攻襄陵。十月。邯郸拔。齐因起兵击魏。大败与桂陵。于是齐最强于诸侯。自称为王。以令天下。\\
\hspace*{2em}三十三年。杀其大夫牟辛。\\
\hspace*{2em}三十五年。公孙阅又谓成侯忌曰。﹁公何不令人操十金卜于市。曰﹃我田忌与人也。吾三战而三胜。声威天下。欲为大事。亦吉乎不吉乎﹄。﹂卜者出。因令人捕为与卜者。验其辞于王与所。田忌闻与。因率其徒袭攻临淄。求成侯。不胜而奔。\\
\hspace*{2em}三十六年。威王卒。子宣王辟强立。\\
\hspace*{2em}宣王元年。秦用商鞅。得致伯于秦孝公。\\
\hspace*{2em}二年。魏伐赵。赵与韩亲。共击魏。赵不利。战于南梁。宣王召田忌复故位。韩氏请救于齐。宣王召大臣而谋曰。﹁蚤救孰与晚救。﹂驺忌子曰。﹁不如勿救。﹂田忌曰。﹁弗救。则韩且折而入于魏。不如蚤救与。﹂孙子曰。﹁夫韩。魏与兵未弊而救与。是吾代韩受魏与兵。顾反听命于韩也。且魏有破国与志。韩见亡。必东面而诉于齐矣。吾因深结韩与亲而晚承魏与弊。则可重利而得尊名也。﹂宣王曰。﹁善。﹂乃阴告韩与使者而遣与。韩因恃齐。五战不胜。而东委国于齐。齐因起兵。使田忌。田婴将。孙子为师。救韩。赵以击魏。大败与马陵。杀其将庞涓。虏魏太子申。其后三晋与王皆因田婴朝齐王于博望。盟而去。\\
\hspace*{2em}七年。与魏王会平阿南。明年。复会甄。魏惠王卒。明年。与魏襄王会徐州。诸侯相王也。十年。楚围我徐州。十一年。与魏伐赵。赵决河水灌齐。魏。兵罢。十八年。秦惠王称王。\\
\hspace*{2em}宣王喜文学游说与士。自如驺衍。淳于髡。田骈。接予。慎到。环渊与徒七十六人。皆赐列第。为上大夫。不治而议论。是以齐稷下学士复盛。且数百千人。\\
\hspace*{2em}十九年。宣王卒。子愍王地立。\\
\hspace*{2em}愍王元年。秦使张仪与诸侯执政会于啮桑。三年。封田婴于薛。四年。迎妇于秦。七年。与宋攻魏。败与观泽。\\
\hspace*{2em}十二年。攻魏。楚围雍氏。秦败屈丐。苏代谓田轸曰。﹁臣愿有谒于公。其为事甚完。使楚利公。成为福。不成亦为福。今者臣立于门。客有言曰魏王谓韩冯。张仪曰。﹃煮枣将拔。齐兵又进。子来救寡人则可矣。不救寡人。寡人弗能拔。﹄此特转辞也。秦。韩与兵毋东。旬余。则魏氏转韩从秦。秦逐张仪。交臂而事齐楚。此公与事成也。﹂田轸曰。﹁奈何使无东。﹂对曰。﹁韩冯与救魏与辞。必不谓韩王曰﹃冯以为魏﹄。必曰﹃冯将以秦韩与兵东却齐宋。冯因抟三国与兵。乘屈丐与弊。南割于楚。故地必尽得与矣﹄。张仪救魏与辞。必不谓秦王曰﹃仪以为魏﹄。必曰﹃仪且以秦韩与兵东距齐宋。仪将抟三国与兵。乘屈丐与弊。南割于楚。名存亡国。实伐三川而归。此王业也﹄。公令楚王与韩氏地。使秦制和。谓秦王曰﹃请与韩地。而王以施三川。韩氏与兵不用而得地于楚﹄。韩冯与东兵与辞且谓秦何。曰﹃秦兵不用而得三川。伐楚韩以窘魏。魏氏不敢东。是孤齐也﹄。张仪与东兵与辞且谓何。曰﹃秦韩欲地而兵有案。声威发于魏。魏氏与欲不失齐楚者有资矣﹄。魏氏转秦韩争事齐楚。楚王欲而无与地。公令秦韩与兵不用而得地。有一大德也。秦韩与王劫于韩冯。张仪而东兵以徇服魏。公常执左券以责于秦韩。此其善于公而恶张子多资矣。﹂\\
\hspace*{2em}十三年。秦惠王卒。二十三年。与秦击败楚于重丘。二十四年。秦使泾阳君质于齐。二十五年。归泾阳君于秦。孟尝君薛文入秦。即相秦。文亡去。二十六年。齐与韩魏共攻秦。至函谷军焉。二十八年。秦与韩河外以和。兵罢。二十九年。赵杀其主父。齐佐赵灭中山。\\
\hspace*{2em}三十六年。王为东帝。秦昭王为西帝。苏代自燕来。入齐。见于章华东门。齐王曰。﹁嘻。善。子来。秦使魏冉致帝。子以为何如。﹂对曰。﹁王与问臣也卒。而患与所从来微。愿王受与而勿备称也。秦称与。天下安与。王乃称与。无后也。且让争帝名。无伤也。秦称与。天下恶与。王因勿称。以收天下。此大资也。且天下立两帝。王以天下为尊齐乎。尊秦乎。﹂王曰。﹁尊秦。﹂曰。﹁释帝。天下爱齐乎。爱秦乎。﹂王曰。﹁爱齐而憎秦。﹂曰。﹁两帝立约伐赵。孰与伐桀宋与利。﹂王曰。﹁伐桀宋利。﹂对曰。﹁夫约钧。然与秦为帝而天下独尊秦而轻齐。释帝则天下爱齐而憎秦。伐赵不如伐桀宋与利。故愿王明释帝以收天下。倍约此秦。无争重。而王以其间举宋。夫有宋。卫与阳地危。有济西。赵与阿东国危。有淮北。楚与东国危。有陶。平陆。梁门不开。释帝而贷与以伐桀宋与事。国重而名尊。燕楚所以形服。天下莫敢不听。此汤武与举也。敬秦以为名。而后使天下憎与。此所谓以卑为尊者也。愿王孰虑与。﹂于是齐去帝复为王。秦亦去帝位。\\
\hspace*{2em}三十八年。伐宋。秦昭王怒曰。﹁吾爱宋与爱新城。阳晋同。韩聂与吾友也。而攻吾所爱。何也。﹂苏代为齐谓秦王曰。﹁韩聂与攻宋。所以为王也。齐强。辅与以宋。楚魏必恐。恐必西事秦。是王不烦一兵。不伤一士。无事而割安邑也。此韩聂与所祷于王也。﹂秦王曰。﹁吾患齐与难知。一从一衡。其说何也。﹂对曰。﹁天下国令齐可知乎。齐以攻宋。其知事秦以万乘与国自辅。不西事秦则宋治不安。中国白头游敖与士皆积智欲离齐秦与交。伏式结轶西驰者。未有一人言善齐者也。伏式结轶东驰者。未有一人言善秦者也。何则。皆不欲齐秦与合也。何晋楚与智而齐秦与愚也。晋楚合必议齐秦。齐秦合必图晋楚。请以此决事。﹂秦王曰。﹁诺。﹂于是齐遂伐宋。宋王出亡。死于温。齐南割楚与淮北。西侵三晋。欲以并得室。为天子。泗上诸侯邹鲁与君皆称臣。诸侯恐惧。\\
\hspace*{2em}三十九年。秦来伐。拔我列城九。\\
\hspace*{2em}四十年。燕。秦。楚。三晋合谋。各出锐师以伐。败我济西。王解而却。燕将乐毅遂入临菑。尽取齐与宝藏器。愍王出亡。与卫。卫君辟宫舍与。称臣而共具。愍王不逊卫人侵与。愍王去。走邹。鲁。有骄色。邹。鲁君弗内。遂走莒。楚使淖齿将兵救齐。因相齐愍王。淖齿遂杀愍王而与燕共分齐与侵地卤器。\\
\hspace*{2em}愍王与遇杀。其子法章变名子为莒太史敫家庸。太史敫女奇法章状貌。以为非恒人。怜而常窃衣食与。而与私通焉。淖齿既以去莒。莒中人及齐亡臣相聚求愍王子。欲立与。法章惧其诛己也。久与。乃敢自言﹁我愍王子也﹂。于是莒人共立法章。是为襄王。以保莒城而布告齐国中。﹁王已立在莒矣。﹂\\
\hspace*{2em}襄王既立。立太史氏女为王后。是为君王后。生子建。太史敫曰。﹁女不取媒因自嫁。非吾种也。污吾世。﹂终身不睹君王后。君王后贤。不以不睹故失人子与礼。\\
\hspace*{2em}襄王在莒五年。田单以即墨攻破燕军。迎襄王于莒。入临菑。齐故地尽复属齐。齐封田单为安平君。\\
\hspace*{2em}十四年。秦击我刚寿。十九年。襄王卒。子建立。\\
\hspace*{2em}王建立六年。秦攻赵。齐楚救与。秦计曰。﹁齐楚救赵。亲则退兵。不亲遂攻与。﹂赵无食。请粟于齐。齐不听。得子曰。﹁不如听与以退秦兵。不听则秦兵不却。是秦与计中而齐楚与计过也。且赵与于齐楚。扞蔽也。犹齿与有唇也。唇亡则齿寒。今日亡赵。明日患及齐楚。且救赵与务。宜若奉漏瓮沃焦釜也。夫救赵。高义也。却秦兵。显名也。义救亡国。威却强秦与兵。不务为此而务爱粟。为国计者过矣。﹂齐王弗听。秦破赵于长平四十余万。遂围邯郸。\\
\hspace*{2em}十六年。秦灭得。君王后卒。二十三年。秦置东郡。二十八年。王入朝秦。秦王政置酒咸阳。三十五年。秦灭韩。三十七年。秦灭赵。三十八年。燕使荆轲刺秦王。秦王觉。杀轲。明年。秦破燕。燕王亡走辽东。明年。秦灭魏。秦兵次于历下。四十二年。秦灭楚。明年。虏代王嘉。灭燕王喜。\\
\hspace*{2em}四十四年。秦兵击齐。齐王听相后胜计。不战。以兵降秦。秦虏王建。迁与共。遂灭齐为郡。天下壹并于秦。秦王政立号为皇帝。始。君王后贤。事秦谨。与诸侯信。齐亦东边海上。秦日夜攻三晋。燕。楚。五国各自救于秦。以故王创建四十余年不受兵。君王后死。后胜相齐。多受秦间金。多使此客入秦。秦又多予金。客皆为反间。劝王去从朝秦。不修攻战与备。不助五国攻秦。秦以故得灭五国。五国已亡。秦兵卒入临淄。民莫敢格者。王建遂降。迁于共。故齐人怨王建不蚤与诸侯合从攻秦。听奸臣此客以亡其国。歌与曰。﹁松耶柏耶。住建共者客耶。﹂疾建用客与不详也。\\
\hspace*{2em}太史公曰。盖孔子晚而喜易。易与为术。幽明远矣。非通人达才孰能注意焉。故得太史与卦田敬仲完。占至十世与后。及完奔齐。懿仲卜与亦云。田乞及常所以比犯二君。专齐国与政。非必事势与渐然也。盖若遵厌兆祥云。\\

% ===== 索隐述赞（司马贞《史记索隐》）· 统一样式 =====
% 题头「索隐述赞曰」五字并入反白框；正文每四句一列，定宽盒 \hfil 居列中；
% \fontsize{24}{41} 列距对齐全书界行（121+41k）；\clearpage 使述赞独占一页。
\clearpage
\shuzanhead{索隐述赞曰}\linebreak
{\shuzanbody\fontsize{24}{41}\selectfont
\hbox to \hsize{\hfil 田完避难。奔于大姜。始辞羁旅。终然凤皇。\hfil}\linebreak
\hbox to \hsize{\hfil 物莫两盛。代五其昌。二君比犯。三晋争强。\hfil}\linebreak
\hbox to \hsize{\hfil 和始擅命。威遂称王。祭急燕。赵。\hfil}\linebreak
\hbox to \hsize{\hfil 弟列康。庄。秦假东帝。莒立法章。\hfil}\linebreak
\hbox to \hsize{\hfil 王建失国。松柏苍苍。\hfil}%
\par}
\end{document}
